\documentclass[11pt]{ltjsarticle}
\usepackage{luatexja-fontspec}
\usepackage{amsmath,amssymb}
\usepackage[margin=22mm]{geometry}
\usepackage{graphicx}
\usepackage{longtable,booktabs,array}
\usepackage{calc}
\usepackage{xcolor}
\usepackage{hyperref}
\hypersetup{colorlinks=true,linkcolor=blue,urlcolor=blue,citecolor=blue}
\makeatletter
\def\maxwidth{\ifdim\Gin@nat@width>\linewidth\linewidth\else\Gin@nat@width\fi}
\makeatother
\setkeys{Gin}{width=\maxwidth,keepaspectratio}
\setcounter{secnumdepth}{0}
\setlength{\parindent}{0pt}
\setlength{\parskip}{0.5em}
\providecommand{\tightlist}{\setlength{\itemsep}{0pt}\setlength{\parskip}{0pt}}
\providecommand{\pandocbounded}[1]{#1}
\providecommand{\real}[1]{#1}
\newcounter{none}
\begin{document}
\section{Appendix: Discreteness Arises Not from an Outer Bound but from
Nodes}\label{appendix-discreteness-arises-not-from-an-outer-bound-but-from-nodes}

\subsection{--- Discreteness as Nodes Generated by
Binarity}\label{discreteness-as-nodes-generated-by-binarity}

\textbf{Author:} Noriaki Kihara (ORCID: 0009-0004-6753-4020)
\textbf{Date:} 18 June 2026 \textbf{DOI (this version):}
10.5281/zenodo.20742991 \textbf{Concept DOI (all versions):}
10.5281/zenodo.20742990 \textbf{Zenodo:}
https://zenodo.org/records/20742991 \textbf{Related paper:} ``On the
Logarithmic Map of Scale and the Separation of Scale and Quantization in
a High-Symmetry Limit,'' §3 (Concept DOI:
\href{https://doi.org/10.5281/zenodo.20740841}{10.5281/zenodo.20740841})

\begin{center}\rule{0.5\linewidth}{0.5pt}\end{center}

\subsection{Purpose}\label{purpose}

This appendix develops independently §3 (``Derivation of discreteness'')
of the related paper. It adds no new claim; its aim is to make explicit
an easily misread point --- that \textbf{discreteness arises
intrinsically not from an outer bound (boundedness) but from the nodes
that Axiom 2 (binarity) generates}.

The explanation ``discreteness comes from boundedness'' implicitly
borrows the eigen-oscillation of a string (only an integer number of
half-waves fit on a string fixed at both ends). But an end (boundary) is
a privileged point with a neighbor only on one side, and violates
anonymity (the lemma of the parent paper). In this system, where ends
are forbidden by anonymity, one cannot borrow the boundary condition of
a string. Stating ``discrete because bounded'' nonetheless smuggles in
the string's boundary condition --- this is circular.

This appendix avoids this circularity and shows that discreteness comes
intrinsically from Axiom 2 (binarity). The root of discreteness is not
the outer bound (boundedness) but the nodes that binarity necessarily
generates.

This appendix is not a rigorous proof but an explication of the
structure of §3 of the parent paper. No physical identification is made.

\begin{center}\rule{0.5\linewidth}{0.5pt}\end{center}

\subsection{1. Axiom 2 generates nodes}\label{axiom-2-generates-nodes}

Axiom 2 of this system (the limit of anonymity = 2) folds the interior
distinction into a binary --- ``no other'' and ``an other present''
(parent paper §0.5, §2). That is, what stands in the interior is binary
(+/−, being/non-being).

When a quantity takes binary values (+/−), to move from the + region to
the − region it must necessarily pass through the boundary of the two.
This boundary --- the point where the sign switches from + to − (or − to
+) --- is called a \textbf{node}.

Being binary itself generates nodes. As long as a + region and a −
region exist, there is necessarily a boundary (a node) between them.
This arises not by imposing a boundary condition from outside, but is
generated intrinsically by binarity (Axiom 2). For a continuous
quantity, a sign switch is merely a single point passed through
smoothly; but under binarity (Axiom 2), the switch from + to − is carved
as a node.

\begin{center}\rule{0.5\linewidth}{0.5pt}\end{center}

\subsection{2. Nodes are discrete}\label{nodes-are-discrete}

A node is a point --- the point of a particular position, the boundary
between + and −. A point is a discrete object, not a continuum.

Therefore, the existence of nodes immediately means the existence of
discrete objects. Between nodes there is a + or − region, and the nodes
themselves line up discretely. Discreteness follows directly from the
existence of nodes as discrete points.

This requires no outer bound (the boundedness of the size of the
system). Nodes arise intrinsically from binarity (Axiom 2), and from the
fact that those nodes are discrete points, discreteness follows. There
is no need to impose an upper bound on the size of the system.

\begin{center}\rule{0.5\linewidth}{0.5pt}\end{center}

\subsection{3. Anonymity binds the nodes to an integer
count}\label{anonymity-binds-the-nodes-to-an-integer-count}

A node is a point (a +/− switch). By anonymity, no node has privilege
--- all nodes are equal (the anonymity of the parent paper).

When equal nodes line up, the number of nodes must be an integer. A node
is binary --- ``present'' or ``not present'' (Axiom 2) --- and a
fractional node (1.5 nodes) cannot exist. Nodes are discrete objects
counted one, two, \ldots, and their total is an integer.

That is, anonymity binds the nodes to an integer count. That the number
of nodes is an integer is a consequence of the nodes being equal
(anonymous) and binary (present/not present).

\begin{center}\rule{0.5\linewidth}{0.5pt}\end{center}

\subsection{4. A closed system binds the nodes to an even
count}\label{a-closed-system-binds-the-nodes-to-an-even-count}

The system is closed (it has no ends and returns upon going around once;
the lemma of the parent paper). When binary (+/−) values alternate on a
closed system, to return to the original sign (+) after one round, the
sign switch must occur an even number of times: + → − → + → \ldots.

Therefore, on a closed system the number of nodes is even. The
requirement of closure --- that the sign returns after one round ---
binds the number of nodes to be even.

Since the interval between nodes corresponds to a half-wavelength, the
number of nodes being even corresponds to the admissible wavelengths
(modes) being numbered by integers (parent paper §3, stage two,
λ\_n=L/n).

\begin{center}\rule{0.5\linewidth}{0.5pt}\end{center}

\subsection{5. The structure of the
derivation}\label{the-structure-of-the-derivation}

From the above, discreteness comes intrinsically not from an outer bound
(boundedness) but from the nodes that binarity (Axiom 2) generates.

{\def\LTcaptype{none} % do not increment counter
\begin{longtable}[]{@{}
  >{\raggedright\arraybackslash}p{(\linewidth - 4\tabcolsep) * \real{0.3333}}
  >{\raggedright\arraybackslash}p{(\linewidth - 4\tabcolsep) * \real{0.3333}}
  >{\raggedright\arraybackslash}p{(\linewidth - 4\tabcolsep) * \real{0.3333}}@{}}
\toprule\noalign{}
\begin{minipage}[b]{\linewidth}\raggedright
Stage
\end{minipage} & \begin{minipage}[b]{\linewidth}\raggedright
Ground
\end{minipage} & \begin{minipage}[b]{\linewidth}\raggedright
Result
\end{minipage} \\
\midrule\noalign{}
\endhead
\bottomrule\noalign{}
\endlastfoot
Genesis of nodes & Axiom 2 (binarity, +/− switch) & nodes arise \\
Discreteness & nodes are points (discrete) & discrete objects exist \\
Integrality & anonymity (nodes equal and binary) & the node count is an
integer \\
Evenness & closure (sign returns after one round) & the node count is
even \\
\end{longtable}
}

The root of discreteness is the nodes that binarity (Axiom 2) generates.
An outer bound (an upper bound on the size of the system) is not needed
for discreteness. On a closed system (the lemma of the parent paper, a
self-unprovable premise), binarity carving the nodes into an integer /
even count makes the admissible wavelengths (modes) numbered by integers
(Figure 1).

\begin{figure}
\centering
\pandocbounded{\includesvg[keepaspectratio]{fig5_nodes_on_closed_loop.pdf}}
\caption{Figure 1: Nodes on a closed loop. Binary (+/−) values
alternate, and their switch points (nodes) arise. Nodes are discrete
points (hence discrete), bound to an integer count by anonymity, and to
an even count by closure (the sign returns to + after one round).
Discreteness comes not from an outer bound but from the nodes that
binarity generates.}
\end{figure}

\textbf{Avoidance of circularity:} ``discrete because bounded'' was a
circularity borrowing the string's boundary condition. The derivation of
this appendix does not borrow a string but derives the nodes
intrinsically from Axiom 2 (binarity). Discreteness is a consequence of
binarity --- an intrinsic structure --- not of the boundedness of the
size of the system.

\begin{center}\rule{0.5\linewidth}{0.5pt}\end{center}

\subsection{6. Limitations (reading
notes)}\label{limitations-reading-notes}

This appendix makes explicit the structure of §3 of the parent paper and
adds no new claim.

\begin{enumerate}
\def\labelenumi{\arabic{enumi}.}
\item
  \textbf{The closed system is a self-unprovable axiom (lemma).} There
  is a line ``anonymity forbids ends → closes'' (the lemma of the parent
  paper), but it is not fully derived from a more fundamental principle.
  The evenness (§4) of this appendix presupposes this closure.
\item
  \textbf{Discreteness comes from Axiom 2 and anonymity.} This appendix
  separates discreteness from the outer bound (boundedness) and
  attributes it to the nodes that Axiom 2 (binarity) generates. However,
  the derivation of Axiom 2 itself is outside the scope of the parent
  paper and this appendix.
\item
  \textbf{``Node'' refers to the switch point of the interior binary.}
  No identification with a physical wave node is made. It is limited to
  the formal structure of the switch of the binary (+/−,
  being/non-being).
\item
  \textbf{Boundedness is unnecessary for discreteness but is separately
  involved in the integrality of mode numbers.} §3, stage two, of the
  parent paper states that the standing wave on a closed system
  (circumference L) is numbered by integer modes λ\_n=L/n.~The evenness
  of nodes (§4) of this appendix corresponds to this. The boundedness
  (the size of the system) itself is not the ground of discreteness, but
  the circumference L of the closed system is involved in the numbering
  of modes.
\item
  \textbf{No physical identification is made.} The physical names of the
  binary, nodes, and modes are anonymous labels, and this appendix makes
  no identification with a particular physical quantity.
\end{enumerate}

\begin{center}\rule{0.5\linewidth}{0.5pt}\end{center}

\subsection{Reference}\label{reference}

\begin{itemize}
\tightlist
\item
  Kihara, N. ``On the Logarithmic Map of Scale and the Separation of
  Scale and Quantization in a High-Symmetry Limit --- A Limited
  Verification Report from a Minimal Axiom System,'' §3, the lemma of §2
  (closure is a consequence of anonymity), Axiom 2. Concept DOI:
  \href{https://doi.org/10.5281/zenodo.20740841}{10.5281/zenodo.20740841}.
\end{itemize}
\end{document}
